\scanpage{149}
de $SL(2,\mathbb Z)$ sur $\mathcal D^+\times\mathbb C$. On peut en fait faire opérer $\mathbb Z^2.\underbrace{SL(2,\mathbb Z).\deleted{\mathbb Z^2}}_{\text{$\frac12$-direct}}$ sur $\mathcal D^+\times\mathbb C$, par
\[
T_{g,m,n}(\tau,z)=\left(g\tau,\frac{z}{c\tau+d}+m+ng\tau\right) \tag*{(15)}
\]
et l'opération déduite de $SL(2,\mathbb Z)$ sur $X_{\mathcal D^+}$ se déduit de (15) en passant au quotient \deleted{\unsure{sur|sous}} dans $\mathcal D^+\times\mathbb C$ par le \deleted{\unsure{g|groupe|faux}}sous-groupe $\mathbb Z^2$ de $\mathbb Z^2.SL(2,\mathbb Z)$.

On peut énoncer dès maintenant

\uline{Théorème} La famille $\mathbb C$-analytique \unsure{de|des} courbes elliptiques, \deleted{\unsure{défini|définí|définie}}\inserted{modulaire} pour la rigidification de Torelli stricte, définie par
\diagram{\begin{tikzcd}[ampersand replacement=\&,row sep=large,column sep=small]
X_{D^+}=D^*\times\mathbb C/\mathbb Z^2 \arrow[d,"{(\tau,\dot z)\mapsto z}"] \& \substack{\mathbb Z^2\ \text{opérant par}\\T_{m,n}(\tau,z)=(\tau,z+m+n\tau)}\\
D^+ \&
\end{tikzcd}}
est un revêtement universel de $M_{1,1}$, avec \deleted{(i.e. \unsure{elle|il} est localement modulaire, et \unsure{épimorphique|épimorphisme},} comme groupe d'automorphismes sur $M_{1,1}$ le groupe $SL(2,\mathbb Z)$, opérant par
\[
g\tau=\frac{a\tau+b}{c\tau+d}\qquad g=\begin{pmatrix}a&b\\c&d\end{pmatrix},\qquad\begin{aligned}a,b,c,d&\in\mathbb Z\\ad-bc&=1\end{aligned}
\]
L'opération précédente est aussi décrite, en termes du pb modulaire de Torelli, par
% archive sheet number: 148
